\documentclass[ ../main.tex]{subfiles}
\providecommand{\mainx}{..}
\begin{document}
\section{Approximate hash map}
\label{sec:pmf}
The approximate map generalizes the Bernoulli set to the \emph{Bernoulli map}~\cite{bmaps}: where the Bernoulli set approximates the indicator function $\mathbf{1}_A \colon U \to \{0,1\}$, the Bernoulli map approximates a general partial function $f \colon X \rightharpoonup Y$.

If the input is not in the domain of definition, the approximate map outputs a random sequence of $d$ bits. If proportion $f$ of the bit patterns represent \emph{undefined}, then the probability that an element not in the domain of definition correctly maps to undefined is $f$ and the probability that it incorrectly maps to some valid value is $1-f$.



The \emph{\gls{gls-pmf}} is a data structure given by the following definition.
\begin{definition}
The \emph{\gls{gls-pmf}} is a data structure that implements the abstract data type of the \emph{static approximate map} where both the \emph{keys} and \emph{values} are from the universe $\BitSet^*$.
A \gls{gls-pmf} that approximates a map with $m$ keys, a \emph{mean} bit length per value $\mu$, false positive rate $\fprate$, and load factor $r$ is denoted by
\begin{equation}
    \pmapf(m, \mu, \fprate, r)\,.
\end{equation}
The \emph{minimal} \gls{gls-pmf} has load factor $r=1$ and may be denoted by $\pmapf(m,\mu,\fprate)$.
\end{definition}

\begin{algorithm}[h]
    \caption{Theoretical implementation of \protect\MakeApproxMap}
    \label{alg:shmap}
    \SetKwProg{func}{function}{}{}
    \KwIn
    {
        $\Mt \in \mathcal{B}^* \times \mathcal{B}^*$ is the \emph{finite map} $\{(k_1,v_1),\ldots,(k_m,v_m)\}$ to approximate and $\fprate$ is the false positive rate.
    }
    \KwOut
    {
        A \gls{gls-pmf} $\Mt^* \supset \Mt$ with a false positive rate $\fprate$.
    }
    \func{\MakeApproxMap{$\Mt$, $\fprate$}}
    {
        \tcp{Make a minimal perfect hash function for the keys.}
        $K \gets \left\{k \colon (k,v) \in \Mt\right\}$\;
        $\ph_{K} \colon \mathcal{B}^* \mapsto \mathbb{Z}^+ \gets$ \MakeMinimalPerfectHash{$K$}\;

        $k \gets \left\lceil \log_2 \frac{1}{\fprate} \right\rceil$\;
        $\matrx{V} \gets m$-by-$k$ bit matrix\;
        \For{$j \gets 1$ \KwTo $m$}
        {
            $\matrx{V}\!\left[\ph_{K}(x_j)\right] \gets \operatorname{h}(k_j) \bmod k$\;
        }
        \For{$n \gets 0$ \KwTo $\infty$}
        {
            $\found \gets \True$\;
            \For{$j \gets 1$ \KwTo $2^n$}
            {
                $b_n \gets$ draw random bit string from $\{0,1\}^n$ without replacement\;
                \For{$i \gets 1$ \KwTo $m$}
                {
                    \tcp{The value $v_j$ is \emph{self-delimiting code} so we
                    check that the first $\operatorname{BL}(v_j)$ bits of the hash is equal to $v_j$.}
                    \If{$\operatorname{h}(k_j \cat b_n) \bmod \operatorname{BL}(v_j^*) \neq v_j^*$}
                    {
                        \label{line:selfterm}
                        $\found \gets \False$\;
                        \Break\;
                    }
                }
                \If{not \found}
                {
                    \Break\;
                }
            }
            \If{\found}
            {
                \tcp{This tuple codes the \gls{gls-pmf}.}
                \Return $(\matrx{V},b_n)$\;
            }
        }
    }
\end{algorithm}


\subsection{Lower-bound on space complexity}
% TODO: Derive the information-theoretic lower bound on the space complexity
% of the approximate map, analogous to Postulate~\ref{pst:approx_l_b} for sets.

\subsection{Upper-bound on entropy}
% TODO: Derive the upper bound on the entropy of the PMF encoding.



\end{document}